Se realizó la reingeniería mecánica parcial y la automatización de un brazo robótico cartesiano para la línea de empaque del Laboratorio de Diseño de Procesos de la Universidad del Valle de Guatemala. El guiado original, formado por tubos metálicos y carros impresos, presentaba fricción, desalineación y avance desigual del pórtico que no se corrigieron con lubricación. En Autodesk Inventor se rediseñaron los carros, el conjunto central y el mecanismo vertical para incorporar ejes de acero de 12~mm y rodamientos lineales, y se fabricaron 48 piezas en PLA.

El control se distribuyó entre una Portenta H7 con Machine Control, que acciona los motores paso a paso y lee los finales de carrera y el encoder de la banda, y una ESP32, que atiende el mando inalámbrico, la pantalla OLED, la pinza y una cámara HUSKYLENS 2. El enlace entre ambos controladores se migró de I2C a UART sobre RS485 por la corrupción persistente de paquetes. La cámara localiza las piezas con un modelo YOLO11n de segmentación y una homografía calibrada con cuatro AprilTags; el encoder permite predecir el instante de captura con la banda en movimiento.

En 120 ciclos automáticos, la pieza se recogió y se entregó en el 61,7~\% de los casos: 73,3~\% con la banda al 25~\% y 50,0~\% al 35~\%. Las piezas verticales alcanzaron 90,6~\%; las horizontales, 45,2~\%, y las diagonales, 54,4~\%. La referenciación mostró una desviación estándar del recorrido medido de 1,3~mm en X y 0,5~mm en Y. No se midieron la exactitud ni la repetibilidad del efector final.
